\subsection*{Simulação CFD}

A simulação foi realizada no OpenFOAM v2606~\cite{openfoam}, segundo o procedimento da
Figura~\ref{fig:diag_cfd}.

\begin{figure}[H]
\centering
\begin{tikzpicture}[node distance=0.45cm]
  \node[etapa] (c1) {\textbf{1.} Geometria e domínio computacional};
  \node[etapa, below=of c1] (c2) {\textbf{2.} Geração da malha};
  \node[etapa, below=of c2] (c3) {\textbf{3.} Condições de fronteira e modelo físico};
  \node[etapa, below=of c3] (c4) {\textbf{4.} Resolução (\texttt{simpleFoam})};
  \node[etapa, below=of c4] (c5) {\textbf{5.} Coeficientes de força};
  \node[etapa, below=of c5] (c6) {\textbf{6.} Pós-processamento};
  \foreach \a/\b in {c1/c2,c2/c3,c3/c4,c4/c5,c5/c6}
    \draw[seta] (\a) -- (\b);
  % ciclo de convergência de malha
  \draw[seta, dashed] (c5.east) -- ++(0.3,0) |- (c2.east);
  % entradas
  \node[io, left=0.5cm of c1, align=flush right] (k1) {Coordenadas do perfil; dimensões do domínio [TODO]};
  \node[io, left=0.5cm of c2, align=flush right] (k2) {Domínio; número de células, $y^+$ pretendido, refinamento junto à parede [TODO]};
  \node[io, left=0.5cm of c3, align=flush right] (k3) {$V_\infty = 100$~m/s, $\nu = 10^{-5}$~m$^2$/s, $c = 1$~m, $\alpha$; modelo de turbulência e intensidade turbulenta [TODO]};
  \node[io, left=0.5cm of c4, align=flush right] (k4) {Caso configurado; esquemas numéricos; critério de paragem dos resíduos [TODO]};
  \node[io, left=0.5cm of c5, align=flush right] (k5) {Campos $p$ e $\mathbf{U}$; $c$, $V_\infty$; direções de sustentação e arrasto};
  \node[io, left=0.5cm of c6, align=flush right] (k6) {Campos convergidos para $\alpha = 0^\circ, 5^\circ, \dots, 20^\circ$};
  % saídas
  \node[io, right=0.75cm of c1] (r1) {Domínio 2D em torno do perfil};
  \node[io, right=0.75cm of c2] (r2) {Malha; qualidade (\texttt{checkMesh}) e $y^+$};
  \node[io, right=0.75cm of c3] (r3) {Caso pronto (\texttt{0/}, \texttt{constant/}, \texttt{system/})};
  \node[io, right=0.75cm of c4] (r4) {Campos $\mathbf{U}$, $p$ e turbulentos; histórico dos resíduos};
  \node[io, right=0.75cm of c5] (r5) {$c_l$, $c_d$ (\texttt{forceCoeffs}); ciclo de refinamento até $c_l$ deixar de variar};
  \node[io, right=0.75cm of c6] (r6) {Linhas de corrente, $C_p$ no perfil, ponto de estagnação, $C_{p,\min}$};
  \foreach \k in {1,...,6} {
    \draw[setaio] (k\k) -- (c\k);
    \draw[setaio] (c\k) -- (r\k);
  }
  \path (c1.north) ++(0,0.6) coordinate (topo);
  \node[titulo] at (k1 |- topo) {Entradas};
  \node[titulo] at (r1 |- topo) {Saídas};
\end{tikzpicture}
\caption{Diagrama de blocos da simulação CFD. A seta a tracejado representa o
estudo de convergência de malha.}
\label{fig:diag_cfd}
\end{figure}

\begin{etapas}
  \item \textbf{Geometria e domínio.} O perfil é importado a partir das
      coordenadas e envolvido por um domínio 2D retangular, de dimensões $x\in[-2;6];\ y\in[-2;2]$, com o bordo de ataque do perfil centrado na origem $(0,0)$.
  %\item \textbf{Malha.} Malha [TODO] com [TODO] células, refinada junto à parede
  %para $y^+ \approx$~[TODO]; a qualidade verifica-se com \texttt{checkMesh}.

  \item \textbf{Malha.} Foi usada uma malha hexaédrica, gerada com o \texttt{snappyHexMesh}, com \approx 65000 células. O primeiro passo foi determinar a altura da primeira \textit{inflation layer}, $h_{FL}$, necessário  para obter $y+\approx1$:
  \begin{equation}
      \mathrm{Re} = \frac{U_\infty \cdot c}{\nu}; \qquad C_f = \frac{0.027}{\mathrm{Re}^{1/7}}; \qquad u_t = U_\infty \cdot \sqrt{\frac{C_f}{2}}; \qquad h_{FL} = 2 \cdot \frac{y^+ \cdot \nu}{u_t}
      \label{eq:y+}
  \end{equation}

  Para $c=1$, $U_\infty=\SI{95}{\meter\per\second}$ e $\nu=1.789\cdot10^{-5}\ \unit{\meter^2\per\second}$, foi obtido $h_{FL}\approx 1\cdot10^{-5}$. De forma a obter uma melhor malha junto à camada limite, foi usada uma $h_{FL} = 1.5\cdot10^{-5}$, que resulta num $y^+$ teórico de $1.5$, pelo que ainda satisfaz a condição de $y^+<5$.

  Foram utilizadas 20 \textit{inflation layers}, com uma taxa de expansão de 1.2, de forma a assegurar uma transição suave entre as \textit{inflation layers} e o restante domínio.

  Foi utilizado uma célula exterior do domínio com $\SI{0.1}{\meter}$ de lado, tendo sido usado uma caixa de refinamento de nível 1 em $x\in[-1;5];\ y\in[-1;1]$. Foi também aplicado um refinamento local, junto à superfície do perfil, com nível de refinamento de 6 (ou 7 para zonas mais críticas).

  A qualidade das malhas foi verificada com o \texttt{checkMesh}.

  %\item \textbf{Condições de fronteira e modelo físico.} Velocidade uniforme
  %$V_\infty$ na entrada, pressão fixa na saída e não escorregamento no perfil;
  %$\alpha$ imposto [TODO]: rodando $\mathbf{V}_\infty$ ou o perfil]. Resolvem-se as
  %equações RANS 2D, estacionárias e incompressíveis, com o modelo de turbulência
  %[TODO], com $Re = V_\infty c/\nu = 10^{7}$, próximo do valor em voo
  %($Re \approx 9{,}6\times10^{6}$, com a corda na raiz e a atmosfera padrão
  %ao nível do mar):
  %\[
  %  \nabla\cdot\overline{\mathbf{U}} = 0, \qquad
  %  (\overline{\mathbf{U}}\cdot\nabla)\,\overline{\mathbf{U}} =
  %  -\nabla \bar{p} + \nabla\cdot\left[(\nu + \nu_t)
  %  \left(\nabla\overline{\mathbf{U}} + \nabla\overline{\mathbf{U}}^{\mathsf{T}}\right)\right],
  %\]
  %onde $\bar{p}$ é a pressão cinemática (dividida por $\rho$).

  \item \textbf{Condições de fronteira e modelo físico.} Para realizar as simulações com diferentes ângulos de ataque (0\textdegree, 5\textdegree, 10\textdegree, 15\textdegree, 20\textdegree), decidiu-se variar o ângulo do perfil, em vez do ângulo do fluxo na entrada. 

  Desta forma, foi utilizada condição de velocidade uniforme $U_\infty = \SI{95}{\meter\per\second}$ na entrada, pressão fixa na saída e não escorregamento no perfil.

  O fluxo aerodinâmico foi modelado recorrendo ao método \textit{Reynolds-averaged Navier-Stokes} (RANS), utilizando como modelo de turbulência o \textit{Shear Stress Transport} (SST) $\mathrm{k-}\omega$, governado pelas seguintes equações:
  \[
    \nabla\cdot\overline{\mathbf{U}} = 0, \qquad
    (\overline{\mathbf{U}}\cdot\nabla)\,\overline{\mathbf{U}} =
    -\nabla \bar{p} + \nabla\cdot\left[(\nu + \nu_t)
    \left(\nabla\overline{\mathbf{U}} + \nabla\overline{\mathbf{U}}^{\mathsf{T}}\right)\right],
  \]
  Este modelo combina as vantagens do $\mathrm{k-}\omega$ junto à camada limite com as vantagens do $\mathrm{k-}\epsilon$ nas regiões distantes da camada limite.

  Foi considerado um $\mathrm{Re}\approx5.31\cdot10^6$, calculado com a \cref{eq:y+}.

  \item \textbf{Resolução.} Algoritmo SIMPLE (\texttt{simpleFoam}) até os
  resíduos descerem abaixo de [TODO] e $c_l$ estabilizar.
  \item \textbf{Coeficientes de força.} Integrando a pressão e as tensões
  viscosas no perfil (\texttt{forceCoeffs}),
  \[
    C_p = \frac{\bar{p} - \bar{p}_\infty}{\tfrac{1}{2}V_\infty^2}, \qquad
    c_l = \frac{L'}{\tfrac{1}{2}\rho_\infty V_\infty^2 c}, \qquad
    c_d = \frac{D'}{\tfrac{1}{2}\rho_\infty V_\infty^2 c}.
  \]
  Repete-se com malhas sucessivamente mais finas até $c_l$ deixar de variar
  (seta a tracejado); a escolha da malha justifica-se nos Resultados.
  \item \textbf{Pós-processamento.} Para $\alpha = 0^\circ, 5^\circ, \dots,
  20^\circ$: linhas de corrente, $C_p$ ao longo do perfil, ponto de estagnação e
  $C_{p,\min}$.
\end{etapas}

